\section{Blaschke geometry and model normalization}
\label{geom:section}

\begin{lemma}\label{geom:cut-time}
Let $D=\inj(M,g)=\diam(M,g)$. Every unit-speed geodesic segment
of length at most $D$ is minimizing. A segment of length less than
$D$ is the unique minimizing segment between its endpoints, and
no segment of length greater than $D$ is minimizing.
\end{lemma}

\begin{proof}
Compactness and positive dimension give $0<D<\infty$, and every
pointwise injectivity radius is at least $D$. For $0<t<D$, the
exponential map at a point $p$ is injective and nonsingular on a
tangent ball containing the sphere of radius $t$. A minimizing
geodesic to $\exp_p(tv)$ exists by completeness and has length at
most $t$. Its initial vector lies in that ball and maps to the same
endpoint as $tv$, so injectivity identifies the two vectors.
This proves both minimality and uniqueness. Taking $t\uparrow D$
and using continuity of distance proves minimality at $D$.
If a segment of length greater than $D$ minimized, the distance
between its endpoints would exceed the diameter.
\end{proof}

In dimension one, a connected closed manifold is a circle.
Arclength identifies its metric with a round circle of the same
length. We henceforth suppose $n=\dim M\ge2$.

We use the following classical structure conclusions
\citep[Chapters~5 and~7]{Besse1978}. Precise statements in the corrected
survey are \citet[Section~4, Corollary~1, Theorem~1, and
Sections~7, 11--12]{McKay2016}.

\begin{proposition}[Classical Blaschke structure]
\label{geom:structure}
Every unit-speed geodesic is simply closed with period $2D$.
For each $p\in M$, the map $v\mapsto\exp_p(Dv)$ from $S_pM$
onto the cut locus is a smooth sphere fibration whose fibers are
great spheres of one dimension $k$, independent of $p$. Its
differential has nullity $k$. The model is determined as follows:
\begin{enumerate}[label=\textup{(\roman*)}]
\item If $k=n-1$, the model is a sphere; if $k=0$, it is a real
projective space. In these two cases the metric classification is
known: the manifold is isometric to its equally scaled model.
\item If $0<k<n-1$, the manifold is simply connected and its
integral cohomology ring has one generator of degree $k+1$.
The possibilities are complex-projective type, quaternionic-projective
type, and Cayley-plane type, with the dimensions listed below.
\end{enumerate}
\end{proposition}

The common period is the prime period, and each geodesic is simple.
Indeed, if two distinct points of its parameter circle represented the same
point of $M$, one of the two arcs between them would have length
at most $D$ and positive length. Lemma~\ref{geom:cut-time} would
make that arc minimizing between equal endpoints, a contradiction.
This also explains the freeness of the period action used later.

The simple-connectivity assertion in the intermediate case also has
a short topological explanation. The homotopy sequence of
$S^k\to S^{n-1}\to\operatorname{Cut}(p)$ gives
$\pi_1(\operatorname{Cut}(p))=0$: the total sphere is simply connected
and the fiber is connected. The exponential map identifies $M$ with
the space obtained by attaching a closed $n$-ball to its cut locus
along this fibration. Indeed, it gives a continuous bijection from
that compact quotient to the Hausdorff manifold $M$, hence a
homeomorphism. Attaching this $n$-cell does not change the fundamental
group, since the intermediate cases have $n\ge4$.

Only the intermediate cases require the new comparison. From now
on rescale the metric so that
\begin{equation}\label{geom:normalization}
 D=\frac\pi2,\qquad T=2D=\pi,\qquad d=n-1,\qquad m=d-k.
\end{equation}
Both $m$ and $k$ are positive. The intermediate model data are
\[
\begin{array}{c|c|c|c}
\text{model}&n&k&m\\ \hline
\CP^r\ (r\ge2)&2r&1&2r-2\\
\HP^r\ (r\ge2)&4r&3&4r-4\\
\OP^2&16&7&8
\end{array}
\]
The projective lines are already included in the sphere case.

On a standard model at this scale, the normal radial curvature
operator has parallel eigenspaces of dimensions $m$ and $k$, with
eigenvalues $1$ and $4$. Solving the scalar Jacobi equations with
initial value zero and derivative one gives the polar density
\begin{equation}\label{geom:model-density}
 w(t)=(\sin t)^m\left|\frac{\sin2t}{2}\right|^k,
 \qquad 0\le t\le T.
\end{equation}
It is positive away from $0,D,T$ and satisfies $w(T-t)=w(t)$.
Let $\sigma_d=|S^d|$ denote ordinary unit-sphere area. The model
volume is
\begin{equation}\label{geom:model-volume}
 V_0=\sigma_d\int_0^D w(t)\,dt.
\end{equation}
The eigenspace description is used only to calculate this model
quantity. No curvature splitting or curvature bound is imposed
on $M$.

Write $U=SM$, let $\phi_t$ be its geodesic flow, and use Liouville
measure with unnormalized angular area:
\begin{equation}\label{geom:liouville}
 d\mu(\xi)=d\operatorname{vol}_g(p)\,d\sigma_p(\xi),
 \qquad \xi\in S_pM.
\end{equation}
This measure is invariant under the flow and has full support;
if $V=\Vol(M,g)$, then $\mu(U)=\sigma_d V$.
Proposition~\ref{geom:structure} gives $\phi_T=\Id_U$.
Below, all determinants of Jacobi tensors are taken in absolute
value, including beyond the cut time.
